% Glossary — shared across chapters (CC.4.4)
\chapter*{Glossary}
\addcontentsline{toc}{chapter}{Glossary}

\begin{description}
  \item[Schema coherence ($\sigma$)] Degree of internal consistency of an agent's representation
    structure; a measurable quantity of the representation space.
  \item[$\sigma$-trap] A bifurcation in the agent's schema coherence: a stable equilibrium at low
    coherence from which gradient-based learning does not escape, associated with compositional
    generalization failure.
  \item[Compositional generalization (CG)] Systematic generalization to novel combinations of
    known components.
  \item[Alignment] The problem of ensuring an AI system's objectives and behavior match human
    values and intent.
  \item[$\Sigma$-Align] The framework developed in this monograph: measure schema coherence,
    intervene on the $\sigma$-trap, and align through coherence.
  \item[Mesa-optimization] The emergence of an inner optimizer or proxy objective that pursues its
    own goal, potentially misaligned with the outer objective.
  \item[Corrigibility] Willingness of a system to allow itself to be corrected or shut down.
  \item[CEV] Coherent Extrapolated Volition — the extrapolation of human values as humans would
    want them to be.
  \item[Indirect Normativity] Using the specification of an idealized normative process rather than
    directly specifying final values.
\end{description}
